%% ─── Introduction ───────────────────────────────────────────────────────────
\section{Introduction}
\label{sec:intro}

Large language models (LLMs) have dramatically changed how developers interact
with codebases, enabling agents to autonomously resolve real-world software
engineering issues~\cite{yang2024sweagent}.  Yet the dominant interaction model remains \emph{ephemeral}:
each agent session begins with an empty context window, forcing the model to
re-discover facts about the codebase from scratch.  This creates a compounding
inefficiency—teams that rely on LLM assistants must accept that every question
triggers a fresh, expensive, and error-prone tour of the filesystem.

Two failure modes characterize this status quo.  First, \textbf{exploration
waste}: agents read far more source files than necessary before reaching an
answer, because there is no prior accumulation of structured knowledge to query.
On representative SWE-Bench tasks, text-based agents spend up to 21 steps
navigating via grep and line-range reads to locate a single target function---work
that structured pre-indexed facts resolve in two or three
queries~\cite{kim2025codestruct}.  Second,
\textbf{context dilution}: irrelevant context degrades LLM
reasoning~\cite{shi2023distracted}, and brittle
string-matching edits fail when whitespace or formatting drifts~\cite{kim2025codestruct}.

Existing retrieval-augmented generation (RAG) systems~\cite{lewis2020rag}
address the first problem partially, but are designed for static document
corpora, not living codebases with interleaved code, documentation, and
decision history.  Code-search benchmarks such as
CodeSearchNet~\cite{husain2019codesearchnet} evaluate individual retrieval
steps in isolation, ignoring the sequential cost of multi-step agent
exploration.

We present \textbf{KB}, a local-first, \emph{facts-first} codebase knowledge
system that:
\begin{enumerate}
  \item[(1)] indexes a repository's code and documentation into a persistent
    SQLite store of atomic, citable facts via deterministic AST analysis and
    sentence-level fact extraction;
  \item[(2)] exposes those facts to LLM agents through a multi-pond BFS
    retrieval loop that merges lexical (BM25), semantic (embedding), and
    graph-walk evidence streams without manual query engineering.
\end{enumerate}

Critically, KB's value is only as compelling as the evidence that it actually
reduces agent exploration cost.  Rubric-based Q\&A scoring—the standard in
prior agent evaluation—captures answer quality but is blind to \emph{how} the
agent arrived at that answer.  A system that reads 40 files before synthesizing
a correct response and one that queries two facts look identical under a
pass/fail quality signal.

To fill this gap, we introduce the \textbf{Multi-Objective Exploration Loss
(MOEL)} evaluation framework.  The harvest runner scores both KB (K) and a
real-agent control (N) on the same eight questions and reports the headline
grade $\Delta S = S_K - S_N$ (Eq.~\ref{eq:delta-s}).  A complementary deep
pipeline runs tasks under three conditions (N, K, O) and aggregates exploration
loss into $L_\text{MOEL}$.

Figure~\ref{fig:system-overview} illustrates the end-to-end system.  Our
contributions are:

\begin{enumerate}
  \item \textbf{KB}: a facts-first codebase knowledge system with deterministic
    AST indexing and a multi-pond hybrid retrieval orchestrator.

  \item \textbf{MOEL}: a multi-objective evaluation framework comprising three
    normalized loss functions ($L_\text{jury}$, $L_\text{trajectory}$,
    $L_\text{resource}$), bias-mitigated LLM jury scoring, and
    logistic-regression calibration on a 20-sample human-annotated set.

  \item \textbf{Empirical results}: build-under-test evaluation compares the
    main and optimized KB binaries against fixed \texttt{kb} and \texttt{raylib}
    repository snapshots.  The optimized build preserves pass rate on both
    targets and keeps fact/document counts stable (Table~\ref{tab:latest-quality}).
    Fresh K-vs-control $\Delta S$ rows are pending because \LatestControlStatus{}.
\end{enumerate}

%% ─── Figure: System Overview placeholder ─────────────────────────────────────
\begin{figure}[t]
\centering
\fbox{
  \parbox{0.92\columnwidth}{
    \centering
    \vspace{0.4cm}
    \textit{\textbf{[FIGURE 1 — System Overview]}} \\[6pt]
    A three-panel horizontal flow diagram: \\
    \textbf{Left panel} — Developer working directory with source files,
      markdown, and git history feeding into \texttt{kb init} (AST indexer +
      sentence segmenter + fact extractor). \\
    \textbf{Center panel} — SQLite knowledge store: \texttt{facts} table
      (doc + code facts), \texttt{fact\_edges} graph, \texttt{fact\_embeddings}
      semantic index, \texttt{documents} (human-readable artifacts). \\
    \textbf{Right panel} — LLM agent loop: \texttt{kb query} dispatches
      multi-pond BFS retrieval, passes ranked facts to LLM, returns grounded
      prose answer. \\
    Below all panels: harvest runner reporting $\Delta S = S_K - S_N$ (K vs.\ real
      agent); deep MOEL harness (N/K/O toy-tool conditions) reporting
      $L_\text{MOEL}$.
    \vspace{0.4cm}
  }
}
\caption{End-to-end KB system.  Primary harvest grade: $\Delta S$ (K vs.\ real
agent); deep MOEL harness adds $L_\text{MOEL}$ under toy-tool N/K/O conditions.}
\label{fig:system-overview}
\end{figure}
